\documentclass[10pt,a4paper]{amsart}
\usepackage[margin=19mm]{geometry}
\usepackage{amsmath,amssymb,mathtools}
\usepackage[scr=rsfs,cal=euler]{mathalfa}
\usepackage{xfrac}
\usepackage[unicode,hidelinks,bookmarks=false]{hyperref}
\newcommand{\ie}{i.e.\ }
\newcommand{\cf}{cf.\ }
\newcommand{\resp}{respectively\ }
\newcommand{\Subsection}{\subsection}
\providecommand{\nobreakdash}{\nobreak\hbox{-}}
\setlength{\emergencystretch}{2em}
\multlinegap0pt
\usepackage{bm}
\usepackage{bbm}
\usepackage{dsfont}
\usepackage{xspace}
\usepackage{cancel}
\usepackage{mathtools}
\usepackage{amsthm}
\makeatletter
\@ifundefined{theorem}{\newtheorem{theorem}{Theorem}}{}
\makeatother
\title[Scalar auxiliary variable (SAV) stabilization of implicit-ex]{Scalar auxiliary variable (SAV) stabilization of implicit-explicit (IMEX) time integration schemes for nonlinear structural dynamics}
\author{Sun-Beom Kwon, Arun Prakash}
\date{Source: arXiv:2406.02839v1 (CC BY 4.0)}
\begin{document}
\maketitle

\noindent\emph{Source and licence.} Scalar auxiliary variable (SAV) stabilization of implicit-explicit (IMEX) time integration schemes for nonlinear structural dynamics. Version arXiv:2406.02839v1 (CC BY 4.0), distributed under the Creative Commons Attribution 4.0 International licence, as recorded in the arXiv licence metadata for this exact version and shown on the abstract page https://arxiv.org/abs/2406.02839v1. Full rights notice: licenses/arxiv-2406.02839-CC-BY-4.0.txt.\par\smallskip

\noindent\textit{Editorial scope.} Counted unit: the Theorem whose source label is \texttt{theorem:4th1} in the article (source0.tex lines 316--326, source label \texttt{theorem:4th1}), delivered together with the article's own complete original proof of it (source0.tex lines 328--388, one \texttt{proof} environment of 61 lines). No mathematics is written, repaired or re-derived: statement and proof are the article's own text line for line. Required context retained uncounted: eqn:4thEnergy (lines 174--177); eqn:4thSAV\_definition (lines 180--183); eqn:4thSAV\_theta (lines 195--198); eqn:4thBDFSAV (lines 219--221); eqn:4thupdated (lines 223--225); eqn:4thupsilon (lines 227--231); eqn:4thxi (lines 227--231). Every definition, display and label of those lines is retained unchanged. Omitted from this package and not redistributed: 1--173, 178--179, 184--194, 199--218, 222--222, 226--226, 232--315, 327--327, 389--1306. Nothing inside a retained excerpt is omitted or altered: every retained body is delivered exactly as the article's lines are written, so no omission receipt of any kind accompanies this package. Citations: the retained excerpts carry no citation command at all, so no bibliography entry of the article is transcribed and no reference is redistributed. Reference adaptation: none was needed and none was made; every reference inside the delivered unit resolves inside it, so no unresolved reference or citation is produced. Printed slips kept literal: no printed word, symbol, hypothesis or number of the counted statement or of its proof was altered. Layout and class: the article's own class is \texttt{elsarticle} with packages including lineno, hyperref, threeparttable, subcaption, amsmath, amssymb, amsthm, algorithm, algcompatible, appendix, caption, comment, diagbox, xcolor; the wrapper is the lane's pinned \texttt{amsart} editorial wrapper at 10pt on A4, which restates the theorem-like environments the retained text needs and adds the article's own macro definitions that the retained text needs (none), with extra wrapper packages (bm, bbm, dsfont, xspace, cancel, mathtools). The delivered numbering is the unit's own. Assets: no retained span contains a figure environment, a graphics-inclusion command, a PDF-page inclusion, a table or a TikZ picture; no image file is copied into this package, the package declares no asset dependency and no image file is redistributed with it. Licence: version arXiv:2406.02839v1 of this article is distributed under the Creative Commons Attribution 4.0 International licence, as recorded in the arXiv licence metadata for this exact version and shown on the abstract page https://arxiv.org/abs/2406.02839v1. A full rights notice is delivered at licenses/arxiv-2406.02839-CC-BY-4.0.txt. Characters: every retained excerpt is pure ASCII, so the delivered engine, XeLaTeX with its default Latin Modern text font, reports no missing character.
\begin{equation}
    \Psi(t) = \frac{1}{2}\dot{\boldsymbol{u}}(t)^T\boldsymbol{M}\dot{\boldsymbol{u}}(t) + \frac{1}{2}\boldsymbol{u}(t)^T\boldsymbol{K}\boldsymbol{u}(t)
    \label{eqn:4thEnergy}
\end{equation}

\begin{equation}
    \Phi(t)=\Psi(t)+\psi
    \label{eqn:4thSAV_definition}
\end{equation}

\begin{equation}
    \Theta=-\dot{\boldsymbol{u}}^T \left( \boldsymbol{M}\ddot{\boldsymbol{u}} + \boldsymbol{K}\boldsymbol{u} \right)=\dot{\boldsymbol{u}}^T \left( \boldsymbol{C}\dot{\boldsymbol{u}}-\boldsymbol{f}^{ext}+\boldsymbol{f}^{NL} \right)
    \label{eqn:4thSAV_theta}
\end{equation}

\begin{align}
    & \frac{\Phi_{n+1}-\Phi_n}{\Delta t} = -\frac{\Phi_{n+1}}{\Psi[\boldsymbol{u}_{n+1}^{IM},\boldsymbol{v}_{n+1}^{IM}]+\psi}\Theta[\boldsymbol{u}_{n+1}^{IM},\boldsymbol{v}_{n+1}^{IM}] \label{eqn:4thBDFSAV}
\end{align}

\begin{align}
    \boldsymbol{u}_{n+1} = \Upsilon^{(k)}_{n+1} \boldsymbol{u}_{n+1}^{IM}, \quad \boldsymbol{v}_{n+1} = \Upsilon^{(k)}_{n+1} \boldsymbol{v}_{n+1}^{IM} \label{eqn:4thupdated}
\end{align}

\begin{align}
    & \Upsilon^{(k)}_{n+1}=1-(1-\Xi_{n+1})^{\beta^{(k)}},
    \quad \textrm{and} \label{eqn:4thupsilon} \\
    & \Xi_{n+1}=\frac{\Phi_{n+1}}{\Psi[\boldsymbol{u}_{n+1}^{IM},\boldsymbol{v}_{n+1}^{IM}]+\psi} \label{eqn:4thxi}
\end{align}

\begin{theorem}
    \label{theorem:4th1}
  For linear structural dynamics, in the absence of external forces $\boldsymbol{f}^{ext}$, the IMEX-BDF$k$-SAV schemes for any $k \ge 1$ are unconditionally energy stable such that
  \begin{equation}
      \Phi_{n+1} \le \Phi_n \label{eqn:4thLTheorem1}
  \end{equation}
  Furthermore, there exists $U^{(k)}>0$ such that
  \begin{equation}
      \Psi[\boldsymbol{u}_n,\boldsymbol{v}_n] \le \left(U^{(k)}\right)^2 \quad  \forall \, n \in [0, T/\Delta t] \label{eqn:4thLTheorem2}
  \end{equation}
\end{theorem}

\begin{proof}
  For linear structural dynamics in the absence of external forces, the definition of $\Theta$ in Eq.~(\ref{eqn:4thSAV_theta}) ensures that $\Theta[\boldsymbol{u}_{n+1}^{IM},\boldsymbol{v}_{n+1}^{IM}] \ge 0$ at all times. 
  Further, from the definition of the SAV $\Phi$ in Eq.~(\ref{eqn:4thSAV_definition}), $\Psi[\boldsymbol{u}_{n+1}^{IM},\boldsymbol{v}_{n+1}^{IM}]+\psi > 0$.
  Thus, given $\Phi_n\ge0$, Eq.~(\ref{eqn:4thBDFSAV}) may be rewritten as:
  \begin{equation}
      \Phi_{n+1} = \dfrac{\Phi_n}{ 1+\Delta t\dfrac{\Theta[\boldsymbol{u}_{n+1}^{IM},\boldsymbol{v}_{n+1}^{IM}]}{\Psi[\boldsymbol{u}_{n+1}^{IM},\boldsymbol{v}_{n+1}^{IM}]+\psi} } 
      \label{eqn:4thTimestepping2}
  \end{equation}  
  ensuring that $\Phi_{n+1} \ge 0$.
  Next, multiplying Eq.~(\ref{eqn:4thBDFSAV}) by $\Delta t$ and noting that the factor on the right-hand side of Eq.~(\ref{eqn:4thBDFSAV}) is simply $\Xi_{n+1}$, as defined in Eq.~(\ref{eqn:4thxi}), we conclude that:
  \begin{equation}
      \Phi_{n+1}-\Phi_n = -\Delta t\Xi_{n+1}\Theta[\boldsymbol{u}_{n+1}^{IM},\boldsymbol{v}_{n+1}^{IM}] \le 0
  \end{equation}
  because, from Eq.~(\ref{eqn:4thxi}), $\Xi_{n+1} \ge 0$.
  This proves Eq.~(\ref{eqn:4thLTheorem1}) of this theorem.
  
  To prove Eq.~(\ref{eqn:4thLTheorem2}), let $U:=\Phi_0$.
  Then, Eq.~(\ref{eqn:4thLTheorem1}) ensures that $\Phi_{n+1} \le U$.
  From Eq.~(\ref{eqn:4thxi}),
  \begin{equation}
      \Xi_{n+1} = \dfrac{\Phi_{n+1}}{\Psi[\boldsymbol{u}_{n+1}^{IM},\boldsymbol{v}_{n+1}^{IM}]+\psi} \le \dfrac{U}{\Psi[\boldsymbol{u}_{n+1}^{IM},\boldsymbol{v}_{n+1}^{IM}]+\psi}
      \label{eqn:4thxi2}
  \end{equation}  
  From Eq.~(\ref{eqn:4thupsilon}), the scalar coefficient $\Upsilon^{(k)}_{n+1}$ can be represented as $\Upsilon^{(k)}_{n+1}=\Xi_{n+1} P_{\beta^{(k)}-1}(\Xi_{n+1})$, where $P_{\beta^{(k)}-1}$ is a polynomial of degree $\beta^{(k)}-1$.
  Note that the absolute value of the polynomial $P_{\beta^{(k)}-1}$ is bounded. 
  Further, since $U>0$ and $\psi>0$, there exists $U^{(k)}>0$ such that 
  \begin{equation}
      |P_{\beta^{(k)}-1}(\Xi^{n+1})| \le \dfrac{2\sqrt{\psi}}{U}U^{(k)}
      \label{eqn:4thPfunction}
  \end{equation} 
  From Eqs.~(\ref{eqn:4thxi2}) and (\ref{eqn:4thPfunction}), 
  \begin{equation}
      |\Upsilon^{(k)}_{n+1}|=|\Xi_{n+1} P_{\beta^{(k)}-1}(\Xi^{n+1})|\le \dfrac{2\sqrt{\psi}U^{(k)}}{\Psi[\boldsymbol{u}_{n+1}^{IM},\boldsymbol{v}_{n+1}^{IM}]+\psi}
      \label{eqn:4thProof1_eta}
  \end{equation}  
  From Eqs.~(\ref{eqn:4thEnergy}) and (\ref{eqn:4thupdated}), the pseudo-energy $\Psi$ at $t_{n+1}$ may be written as:
  \begin{equation}
       \Psi[\boldsymbol{u}_{n+1},\boldsymbol{v}_{n+1}] = \left(\Upsilon^{(k)}_{n+1}\right)^2 \Psi[\boldsymbol{u}_{n+1}^{IM},\boldsymbol{v}_{n+1}^{IM}] 
       \label{eqn:4thPsinextstep}
  \end{equation}
  Substituting Eq.~(\ref{eqn:4thProof1_eta}) into Eq.~(\ref{eqn:4thPsinextstep}), the following inequality is obtained:
  \begin{equation}
      \begin{array}{ll}
           \Psi[\boldsymbol{u}_{n+1},\boldsymbol{v}_{n+1}] 
           &\le \left(\dfrac{2\sqrt{\psi}U^{(k)}}{\Psi[\boldsymbol{u}_{n+1}^{IM},\boldsymbol{v}_{n+1}^{IM}]+\psi}\right)^2 \Psi[\boldsymbol{u}_{n+1}^{IM},\boldsymbol{v}_{n+1}^{IM}] \\
           &\le \left(U^{(k)}\right)^2
      \end{array}        
  \end{equation}
    Note that we use the following identity in the last step above:
    \begin{equation}
      \begin{array}{ll}
            &\dfrac{4 \psi \Psi[\boldsymbol{u}_{n+1}^{IM},\boldsymbol{v}_{n+1}^{IM}] }{\left(\Psi[\boldsymbol{u}_{n+1}^{IM},\boldsymbol{v}_{n+1}^{IM}]+\psi\right)^2} \le 1 \\
      \end{array} 
      \label{eqn:4thIdentity1}
    \end{equation}
    which in turn is obtained from the following statement:
    \begin{equation}
        {\left(\Psi[\boldsymbol{u}_{n+1}^{IM},\boldsymbol{v}_{n+1}^{IM}]-\psi\right)^2} = {\left(\Psi[\boldsymbol{u}_{n+1}^{IM},\boldsymbol{v}_{n+1}^{IM}]+\psi\right)^2} - 4 \psi \Psi[\boldsymbol{u}_{n+1}^{IM},\boldsymbol{v}_{n+1}^{IM}] \ge 0
    \end{equation}
  The proof is complete.
\end{proof}

\end{document}
